\subsection{Algebras of measures}

Let G be a locally compact group. It will be understood, once and for all, that the measures \(\mu_{1},\ldots,\mu_{n}\) on G are said to be convolvable if they are so for the mapping

\[
(x _ {1}, x _ {2}, \ldots , x _ {n}) \mapsto x _ {1} x _ {2} \cdot \cdot \cdot x _ {n};
\]

and it is by means of this mapping that the convolution product \(\ast_{i}\mu_{i}\) will always be taken. If \(s\in G, t\in G\) , then

\[
\varepsilon_ {s} * \varepsilon_ {t} = \varepsilon_ {s t}.\tag{1}
\]

If \(s \in \mathbf{G}\) and \(\mu \in \mathcal{M}(\mathbf{G})\) , then

\begin{enumerate}
\def\labelenumi{(\arabic{enumi})}
\setcounter{enumi}{1}
\tightlist
\item
\end{enumerate}

\[
\begin{array}{l} \varepsilon_ {s} * \mu = \gamma (s) \mu \\ \mu * \varepsilon_ {s} = \delta (s ^ {- 1}) \mu \end{array}\tag{3}
\]

by §1, No.~1, Example 3. If G is abelian, to say that \(\mu_{1}\) and \(\mu_{2}\) are convolvable is equivalent to saying that \(\mu_{2}\) and \(\mu_{1}\) are convolvable, and one then has \(\mu_{1} * \mu_{2} = \mu_{2} * \mu_{1}\) . When G is not abelian, it can happen that \(\mu_{1}\) and \(\mu_{2}\) are convolvable, without \(\mu_{2}\) and \(\mu_{1}\) being so (Exer. 12).

PROPOSITION 1. --- Let G be a locally compact group, \(\lambda, \mu, \nu\) measures \(\neq 0\) on G.

\begin{enumerate}
\def\labelenumi{(\roman{enumi})}
\tightlist
\item
  If \(\lambda, \mu, \nu\) are convolvable, then so are \(\lambda\) and \(\mu\) , \(|\lambda| * |\mu|\) and \(\nu\) , \(\mu\) and \(\nu\) , \(\lambda\) and \(|\mu| * |\nu|\) , and one has
\end{enumerate}

\[
\lambda * \mu * \nu = (\lambda * \mu) * \nu = \lambda * (\mu * \nu).
\]

\begin{enumerate}
\def\labelenumi{(\roman{enumi})}
\setcounter{enumi}{1}
\tightlist
\item
  If \(\lambda\) and \(\mu\) are convolvable, as well as \(|\lambda| * |\mu|\) and \(\nu\) , then \(\lambda, \mu, \nu\) are convolvable. Similarly if \(\mu\) and \(\nu\) are convolvable, as well as \(\lambda\) and \(|\mu| * |\nu|\) .
\end{enumerate}

This follows from Prop. 1 of §1, No.~2.

\pandocbounded{\includegraphics[keepaspectratio]{images/part001_0fdec0aa.jpg}}

There can exist measures \(\lambda, \mu, \nu\) on \(G\) such that the convolution products \(\lambda * \mu\) , \((\lambda * \mu) * \nu\) , \(\mu * \nu\) , \(\lambda * (\mu * \nu)\) are all defined, and yet \((\lambda * \mu) * \nu \neq \lambda * (\mu * \nu)\) (cf.~Exer. 4).

Let \(\rho\) be a lower semi-continuous finite function \(>0\) on \(\mathbf{G}\) such that \(\rho(st) \leqslant \rho(s)\rho(t)\) for all \(s, t\) in \(\mathbf{G}\) . We denote by \(\mathcal{M}^{\rho}(\mathbf{G})\) the vector space of measures \(\lambda\) on \(\mathbf{G}\) such that \(\rho\) is \(\lambda\) -integrable, and by \(\|\lambda\|_{\rho}\) (or simply \(\|\lambda\|\) ) the norm \(\int_{\mathbf{G}} \rho(s)d|\lambda|(s)\) on this space. When \(\rho = 1\) , one recovers the set \(\mathcal{M}^{1}(\mathbf{G})\) of bounded measures on \(\mathbf{G}\) .

PROPOSITION 2. --- (i) Any two elements of \(\mathcal{M}^{\rho}(\mathrm{G})\) are convolvable.

\begin{enumerate}
\def\labelenumi{(\roman{enumi})}
\setcounter{enumi}{1}
\item
  For convolution, and for the norm \(\| \lambda \|\) , \(\mathcal{M}^{\rho}(\mathbf{G})\) is a complete normed algebra, admitting \(\varepsilon_{e}\) as unity element.
\item
  \(\mathcal{C}'(\mathbf{G})\) is a subalgebra of \(\mathcal{M}^{\rho}(\mathbf{G})\) .
\end{enumerate}

Let \(\lambda, \mu\) be in \(\mathcal{M}^{\rho}(\mathbf{G})\) , and let us show that \(\lambda\) and \(\mu\) are convolvable. Let \(f \in \mathcal{K}_{+}(\mathbf{G})\) . Since \(\rho\) is \(>0\) and lower semi-continuous, there exists a constant \(k > 0\) such that \(f \leqslant k\rho\) . Then

\[
\begin{array}{r l} & {\int^ {*} f (s t) d | \lambda | (s) d | \mu | (t) \leqslant k \int^ {*} \rho (s t) d | \lambda | (s) d | \mu | (t)} \\ & {\qquad \leqslant k \int^ {*} \rho (s) \rho (t) d | \lambda | (s) d | \mu | (t)} \\ & {\qquad = k \bigg (\int^ {*} \rho (s) d | \lambda | (s) \bigg) \bigg (\int^ {*} \rho (t) d | \mu | (t) \bigg)} \end{array}
\]

(Ch. V, §8, No.~3, Cor. 1 of Prop. 8). Therefore \((s,t) \mapsto f(st)\) is \((\lambda \otimes \mu)\) -integrable, so that \(\lambda\) and \(\mu\) are convolvable. On the other hand, using Ch. V (§1, Prop. 4, §6, Prop. 2, §8, Cor. 1 of Prop. 8) and the fact that \((s,t) \mapsto \rho(s)\rho(t)\) is lower semi-continuous in \(G \times G\) , one has

\[
\begin{array}{r l} & {\int_ {\mathrm{G}} ^ {*} \rho (s) d | \lambda * \mu | (s) = \int_ {\mathrm{G}} ^ {\bullet} \rho (s) d | \lambda * \mu | (s)} \\ & {\quad \leqslant \int_ {\mathrm{G} \times \mathrm{G}} ^ {\bullet} \rho (s t) d | \lambda | (s) d | \mu | (t) \leqslant \int_ {\mathrm{G} \times \mathrm{G}} ^ {\bullet} \rho (s) \rho (t) d | \lambda | (s) d | \mu | (t)} \\ & {\quad = \int_ {\mathrm{G} \times \mathrm{G}} ^ {*} \rho (s) \rho (t) d | \lambda | (s) d | \mu | (t) = \| \lambda \| \cdot \| \mu \|.} \end{array}
\]

One sees that \(\lambda * \mu \in \mathcal{M}^{\rho}(\mathrm{G})\) and that \(\| \lambda * \mu \| \leqslant \| \lambda \| \cdot \| \mu \|\) . In view of Prop. 1, \(\mathcal{M}^{\rho}(\mathrm{G})\) is an algebra. The mapping \(\lambda \mapsto \rho \cdot \lambda\) is an isometric linear mapping \(\theta\) of \(\mathcal{M}^{\rho}(\mathrm{G})\) into \(\mathcal{M}^{1}(\mathrm{G})\) ; if \(\mu \in \mathcal{M}^{1}(\mathrm{G})\) then \(1 / \rho\) , which is locally bounded and upper semi-continuous, is locally \(\mu\) -integrable, and \(\rho\) is \((1 / \rho) \cdot \mu\) -integrable, thus \((1 / \rho) \cdot \mu \in \mathcal{M}^{\rho}(\mathrm{G})\) ; this proves that \(\theta\) is surjective; therefore \(\mathcal{M}^{\rho}(\mathrm{G})\) is a complete normed algebra. Finally, it is clear that \(\varepsilon_{e}\) is a unity element for \(\mathcal{M}^{\rho}(\mathrm{G})\) and that \(\mathcal{C}'(\mathrm{G})\) is a subalgebra of \(\mathcal{M}^{\rho}(\mathrm{G})\) (§1, No.~4, Cor. of Prop. 5).

If \(\rho = 1\) , Prop. 2, (i) and (ii) also follow from §1, Prop. 2.

PROPOSITION 3. --- Let \(\mu_1, \ldots, \mu_n\) be measures on G. If all of the \(\mu_i\) , except at most one, have compact support, then the \(\mu_i\) are convolvable.

For, let \(S_{i}\) be the support of \(\mu_{i}\) , and suppose that \(S_{i}\) is compact for \(i \neq i_{0}\) . Let K be a compact subset of G. The set of \((x_{1}, \ldots, x_{n}) \in \prod_{i} S_{i}\) such that \(x_{1}x_{2} \cdots x_{n} \in K\) is compact, because the conditions \(x_{i} \in S_{i}\) for all \(i, x_{1}x_{2} \cdots x_{n} \in K\) imply

\[
x _ {i _ {0}} \in \mathrm{S} _ {i _ {0} - 1} ^ {- 1} \dots \mathrm{S} _ {1} ^ {- 1} \mathrm{KS} _ {n} ^ {- 1} \dots \mathrm{S} _ {i _ {0} + 1} ^ {- 1}.
\]

Therefore the \(\mu_{i}\) are convolvable (§1, No.~4, Prop. 4).

PROPOSITION 4. --- The mapping \((\lambda, \mu) \mapsto \lambda * \mu\) (resp. \((\lambda, \mu) \mapsto \mu * \lambda)\) , where \(\lambda \in \mathcal{C}'(G)\) , \(\mu \in \mathcal{M}(G)\) , defines on \(\mathcal{M}(G)\) the structure of a left (resp. right) module over the algebra \(\mathcal{C}'(G)\) .

This follows from Props. 1 and 3.

PROPOSITION 5. --- Let \(\lambda\) be a left (resp. right) Haar measure on G, and \(\mu \in \mathcal{M}^1(\mathrm{G})\) . Then \(\mu\) and \(\lambda\) (resp. \(\lambda\) and \(\mu\) ) are convolvable, and \(\mu * \lambda = \mu(1)\lambda\) (resp. \(\lambda * \mu = \mu(1)\lambda\) ).

We can suppose that \(\mu \geqslant 0\) . Let \(f \in \mathcal{K}_{+}(\mathbf{G})\) . When \(\lambda\) is a left Haar measure,

\[
\int^ {*} d \mu (x) \int^ {*} f (x y) d \lambda (y) = \int^ {*} d \mu (x) \int f (y) d \lambda (y) = \lambda (f) \| \mu \|,
\]

therefore the function \((x,y)\mapsto f(xy)\) is \((\mu \otimes \lambda)\) -integrable, and its integral for \(\mu \otimes \lambda\) is \(\lambda (f)\| \mu \|\) . One argues similarly when \(\lambda\) is a right Haar measure.

PROPOSITION 6. --- Let \(\mu\) and \(\nu\) be two convolvable measures on \(\mathbf{G}\) . Let \(\chi\) be a continuous representation of \(\mathbf{G}\) in \(\mathbf{C}^*\) . Then \(\chi \cdot \mu\) and \(\chi \cdot \nu\) are convolvable and \((\chi \cdot \mu) * (\chi \cdot \nu) = \chi \cdot (\mu * \nu)\) .

Let \(f \in \mathcal{K}(\mathbf{G})\) . Then \(f\chi \in \mathcal{K}(\mathbf{G})\) , therefore the function

\[
(x, y) \mapsto f (x y) \chi (x y) = f (x y) \chi (x) \chi (y)
\]

on \(\mathbf{G} \times \mathbf{G}\) is integrable for \(\mu \otimes \nu\) ; therefore the function \((x, y) \mapsto f(xy)\) is integrable for \((\chi \cdot \mu) \otimes (\chi \cdot \nu)\) ; therefore \(\chi \cdot \mu\) and \(\chi \cdot \nu\) are convolvable. Moreover,

\[
\begin{array}{r l} \langle \chi \cdot \mu * \chi \cdot \nu , f \rangle & = \int f (x y) \chi (x) \chi (y) d \mu (x) d \nu (y) \\ & = \int (f \chi) (x y) d \mu (x) d \nu (y) = \langle \mu * \nu , \chi f \rangle , \end{array}
\]

whence \((\chi \cdot \mu) * (\chi \cdot \nu) = \chi \cdot (\mu * \nu)\) .

PROPOSITION 7. --- Let G and G' be two locally compact groups, u a continuous representation of G in G', and \(\mu_{1},\ldots,\mu_{n}\) measures on G, all \(\neq0\) . Then the following assertions are equivalent:

\begin{enumerate}
\def\labelenumi{(\roman{enumi})}
\item
  \(u\) is \(\mu_{i}\) -proper for all \(i\) , and the measures \(u(|\mu_i|)\) are convolvable;
\item
  the \(\mu_{i}\) are convolvable and \(u\) is proper for \(*_{i}(|\mu_{i}|)\)
\end{enumerate}

When these conditions are satisfied,

\[
\underset {i} {\ast} u (\mu_ {i}) = u \Big (\underset {i} {\ast} \mu_ {i} \Big).
\]

This follows from §1, No.~2, Cor. of Prop. 1.

COROLLARY. --- Let G be a locally compact group, \(\mu_{1},\ldots,\mu_{n}\) measures on G. For the sequence \((\mu_{i})_{1\leqslant i\leqslant n}\) to be convolvable, it is necessary and sufficient that the sequence \((\breve{\mu}_{n-i})_{0\leqslant i\leqslant n-1}\) be so, in which case

\[
\left(\mu_ {1} * \dots * \mu_ {n}\right) ^ {\vee} = \stackrel {\vee} {\mu} _ {n} * \dots * \stackrel {\vee} {\mu} _ {1}.
\]

This follows from Prop. 7 on considering the isomorphism \(x \mapsto x^{-1}\) of \(G\) onto the opposite group.

